\label{sec:methodology}

\subsection{How \textsc{bisect} Handles Prison Population Adjustment}

The \textsc{bisect} pipeline treats prison population adjustment as a
data preprocessing step, implemented in the \texttt{bisect-data} crate
and triggered by the \texttt{--adjust-prison-population} flag on
\texttt{bisect fetch}. This design keeps the adjustment entirely outside
the redistricting algorithm itself: \textsc{bisect}'s partitioning engine
operates on whatever population data is provided, and the adjustment
determines what that data contains. The algorithm is neutral with respect
to the adjustment choice; the data layer implements whatever state law requires.

\begin{verbatim}
# Download unadjusted Census data (default):
bisect fetch --year 2020 --workers 8

# Download with prison population adjustment
# (applied where state law requires):
bisect fetch --year 2020 --workers 8 --adjust-prison-population

# Build plan on adjusted data:
bisect build official_2020 --year 2020 --workers 8
\end{verbatim}

\subsection{Adjustment Algorithm}

The adjustment proceeds in three steps.

\textbf{Step 1: Identify group quarters population.} The Census Bureau's
2020 DHC-A file provides population counts for each type of group quarters
(GQ) facility, including correctional facilities for adults, at the census
tract level. The \texttt{bisect-data} module reads GQ Type 101 (Federal
correctional facilities) and GQ Type 201 (State correctional facilities)
to identify all tracts containing prisons.

\textbf{Step 2: Subtract prison population from tract.} For each
prison-containing tract, the incarcerated population (GQ Types 101 and 201)
is subtracted from the tract's total population. The tract's minority
population composition is adjusted proportionally using the DHC-A
demographic breakdown by GQ type.

\textbf{Step 3: Reallocate to home county.} The subtracted population is
redistributed to the home counties of incarcerated people using the BJS
National Prisoner Statistics (NPS) program data, which reports the
home-state distribution of state prisoners and the home-county distribution
for states that participate in the NPS county-level reporting. For states
without county-level NPS data, the adjustment uses the state-level
distribution weighted by county population.

\subsection{Data Invariants}

The prison population adjustment must satisfy two invariants, which are
validated as L0 tests in the \textsc{bisect} test suite:

\begin{enumerate}
\item \textbf{Population conservation}: The total state population is
  identical before and after adjustment. Prison adjustment reallocates
  people between geographic units; it does not add or remove anyone from
  the count. Formally: $\sum_i p_i^{\text{adj}} = \sum_i p_i^{\text{raw}}$
  for all census tracts $i$ in the state.

\item \textbf{Non-negativity}: No census tract has negative population
  after adjustment. In practice this is always satisfied because prisons
  do not contain more people than the tract's total population---the GQ
  population is a subset of the total---but the invariant is checked
  explicitly: $p_i^{\text{adj}} \geq 0$ for all $i$.
\end{enumerate}

\subsection{North Carolina Case Study}

North Carolina illustrates the adjustment's mechanics. The state has
multiple correctional facilities, but two are particularly relevant for
redistricting: Hoke Correctional Institution in Hoke County and Lumberton
Correctional Institution in Robeson County.

Hoke County, which contains Hoke Correctional Institution, has a total
2020 Census population of approximately 53,600. The facility holds
approximately 900 inmates. Under the unadjusted count, all 900 are counted
in Hoke County. Under the adjustment, these 900 individuals are reallocated
to their home counties---primarily Mecklenburg County (Charlotte),
Durham County, and Forsyth County (Winston-Salem), which have the highest
incarcerated-population origination rates in the state as reported in
North Carolina Department of Adult Correction administrative data.

This reallocation shifts approximately 900 people out of the 11th or
9th congressional district (depending on which encompasses Hoke County
in the plan year) and into the congressional districts covering Mecklenburg
and Durham counties, which are the 12th and 4th districts respectively
in the 2022 congressional plan.

The shift is small---900 people against a congressional district target
of approximately 745,000 in North Carolina---but it is not negligible in
borderline cases: a district that is 1,500 people over quota can be brought
into compliance by a 1,500-person reallocation without any boundary change.
More importantly, the minority composition of the reallocated population
affects the VRA analysis, as discussed in Section~\ref{sec:results}.
